\documentclass{article}
\usepackage{amsmath}
\usepackage{amssymb}
\usepackage{braket}

\title{Mathematical Frameworks of Quantum Mechanics in Math Explorer}
\author{Math Explorer Team}
\date{\today}

\begin{document}

\maketitle

\section{Introduction}
This document outlines the implementation of key mathematical frameworks used in quantum mechanics within the \texttt{math\_explorer} library. The implementation focuses on Linear Algebra, Calculus, Probability Theory, Complex Analysis, and Group Theory, forming a rigorous foundation for quantum simulations.

\section{Linear Algebra: The Skeleton of QM}
The library models quantum states and operators using the \texttt{nalgebra} crate, extended with complex number support.

\subsection{Hilbert Spaces}
A Hilbert space $\mathcal{H}$ is an abstract vector space with an inner product. In our implementation, a quantum state $|\psi\rangle$ is represented by the \texttt{QuantumState} struct, which wraps a complex vector.

\begin{itemize}
    \item \textbf{State Vectors}: $|\psi\rangle$ is a column vector.
    \item \textbf{Dual Vectors}: $\langle\psi|$ is the conjugate transpose.
    \item \textbf{Inner Product}: $\langle\phi|\psi\rangle$ is implemented via the \texttt{inner\_product} method.
    \item \textbf{Normalization}: States are normalized such that $\langle\psi|\psi\rangle = 1$.
\end{itemize}

\subsection{Linear Operators}
Physical observables are represented by the \texttt{QuantumOperator} struct.
\begin{itemize}
    \item \textbf{Hermitian Operators}: Operators $\hat{A}$ such that $\hat{A} = \hat{A}^\dagger$ correspond to observables.
    \item \textbf{Expectation Values}: The average measurement result is given by $\langle\hat{A}\rangle = \langle\psi|\hat{A}|\psi\rangle$.
\end{itemize}

\section{Calculus and Dynamics}
The time evolution of a quantum system is governed by the Schrödinger Equation:
\begin{equation}
    i\hbar \frac{\partial}{\partial t}|\psi(t)\rangle = \hat{H}|\psi(t)\rangle
\end{equation}
For a time-independent Hamiltonian $\hat{H}$, the formal solution is:
\begin{equation}
    |\psi(t)\rangle = e^{-i\hat{H}t/\hbar}|\psi(0)\rangle
\end{equation}
This is implemented in the \texttt{schrodinger} module using matrix exponentiation.

\section{Probability Theory}
According to the Born Rule, the probability of finding a system in a basis state is the square of the magnitude of the projection.
\begin{equation}
    P(x) = |\psi(x)|^2
\end{equation}
The \texttt{probability\_density} method returns these probabilities.

\section{Group Theory and Spin}
The library includes implementations of the Pauli matrices, which are generators of the SU(2) group, essential for describing spin-$\frac{1}{2}$ systems.
\begin{equation}
    \sigma_x = \begin{pmatrix} 0 & 1 \\ 1 & 0 \end{pmatrix}, \quad
    \sigma_y = \begin{pmatrix} 0 & -i \\ i & 0 \end{pmatrix}, \quad
    \sigma_z = \begin{pmatrix} 1 & 0 \\ 0 & -1 \end{pmatrix}
\end{equation}
Commutators $[\hat{A}, \hat{B}]$ quantify the simultaneous measurability of observables.

\section{Fourier Analysis}
To switch between position and momentum representations in discrete systems, the Discrete Fourier Transform (DFT) operator is provided.

\end{document}
